% ================================================================
% SECTION 4 — THE FIVE ARCHITECTURAL PILLARS OF BX3
% Expanded treatment — ~450 lines (vs prior 320)
% ================================================================
\section{The Five Architectural Pillars}\label{sec:pillars}

The \bxthree{} upstream accountability guarantee is realized through five
named architectural pillars.  Each pillar addresses a distinct failure mode
in multi-agent systems and is independently formalizable and auditable.  Together
they constitute the full structural commitment: that \bxthree{} systems fail upward
into human consciousness, never downward into autonomous chaos.

The five pillars are:

\begin{enumerate}\itemsep=0pt
\item[\textbf{P1 — Loop Isolation.}] Prevents state circularity in recursive
execution graphs; ensures provenance of every decision is traceable to its
originating node.

\item[\textbf{P2 — Recursive Spawning.}] Maintains immutable parent-child
accountability chains as nodes spawn subordinate agents; provides the
propagation substrate for the Bailout Protocol.

\item[\textbf{P3 — Spatial Firewall.}] Enforces execution domain separation
between spawned sub-agents; prevents cross-domain state contamination.

\item[\textbf{P4 — Sandbox Gate.}] Provides controlled inter-domain
communication with mandatory pre-commit hooks; makes Bounds Engine
operating range enforcement structurally operative.

\item[\textbf{P5 — Bailout Protocol.}] Compels every unresolved exception
to propagate to the human anchor $h(n)$ via the parent chain, bypassing all
intermediate machine actors with immutable forensic attribution.
\end{enumerate}

Each pillar is specified in its own dedicated section below, with formal
properties, enforcement mechanisms, relationship to the three-layer model,
implementation details, and contribution to the upstream accountability
guarantee.